\documentclass[11pt,a4paper]{article}

% --- Pacchetti Matematici Fondamentali ---
\usepackage{amsmath}   % Strutture matematiche avanzate (es. align)
\usepackage{amssymb}   % Simboli matematici e font speciali (es. \mathbb pour i reali R)
\usepackage{amsfonts}  % Font matematici aggiuntivi
\usepackage{amsthm}    % Gestione di teoremi, definizioni e dimostrazioni

% --- Configurazione dei Teoremi ---
\newtheorem{theorem}{Teorema}
\newtheorem{lemma}[theorem]{Lemma}
\newtheorem{proposition}[theorem]{Proposizione}
\theoremstyle{definition}
\newtheorem{definition}[theorem]{Definizione}
\newtheorem{example}[theorem]{Esempio}
\theoremstyle{remark}
\newtheorem{remark}[theorem]{Nota}

\begin{document}

\begin{theorem}

	Let $p,M,c$ such that $\Delta-tp(c) = \{ \phi \in \Delta : p \vdash \phi\}$. Then $p$ prime.
 
\end{theorem}

\begin{proof}

	By contradiction, assume there exist $\phi_1,..,\phi_n \in P(\Delta)$ such that $p \vdash \bigvee_{i=1,..,n} \phi_i$ 
	but for all $i=1,..,n$ we have $p \not \vdash \phi_i$. Since $c \models p$ we must have $c \models \phi_i$ for some $i = 1,..,n$.
	Put $\phi_i$ in $DNF$ to obtain $c \models \bigvee_{k=1,..,m} \bigwedge_{j=1,..,n_k} \psi_j^k$ where $\psi_j^k \in \Delta$ for all $j,k$.
	Let $k$ such that $c \models \bigwedge_{j=1,..,n_k} \psi_j^k$.
	Then for all $j = 1,..,n_k$ we have $c \models \psi_j^k$. 
	Then for all $j=1,..,n_k$ we have $p \vdash \psi_j^k$ and consequently $p \vdash \phi_i$, contradiction.

\end{proof}

\end{document}

